% !Mode:: "TeX:UTF-8"

\chapter{实验结果与分析框架}

本章为后续完善报告时的结果承载框架。若真实模型采集已经完成，可直接从 \texttt{outputs/tables} 与 \texttt{outputs/figures} 中补充统计结果。

\section{整体准确率}

建议在此处插入 \texttt{outputs/tables/accuracy\_by\_level.csv} 的关键数值，并结合图 \ref{fig:accuracy} 说明不同模型在 L1、L2、L3 下的准确率差异。

\begin{figure}[htbp]
    \centering
    \includegraphics[width=0.82\textwidth]{accuracy_heatmap.png}
    \bicaption{不同模型在 Prompt 层级上的准确率热力图}{Accuracy heatmap across prompt levels}
    \label{fig:accuracy}
\end{figure}

\section{模型行为边界}

图 \ref{fig:radar} 可用于对比不同模型在准确率、一致性、公式使用、坐标建模和噪声敏感性等维度上的综合表现。正式写作时，应结合具体指标解释各模型的优势和边界。

\begin{figure}[htbp]
    \centering
    \includegraphics[width=0.78\textwidth]{model_radar.png}
    \bicaption{模型行为特征雷达图}{Radar chart of model behavioral features}
    \label{fig:radar}
\end{figure}

\section{无关噪声对判断的影响}

该部分可结合 \texttt{outputs/tables/noise\_flip\_rate.csv}，分析不同噪声条件是否引起模型判断翻转。重点关注 borderline 场景，因为临界场景更容易放大模型的不稳定性。

\section{失效模式分析}

可结合 \texttt{outputs/tables/failure\_modes.csv} 和错误样例，分析错误主要集中在哪些类别。建议至少选取 2 到 3 条典型回复，说明模型如何出现概念混淆、计算崩溃或噪声牵引。

\section{关联规则发现}

可从 \texttt{outputs/tables/outcome\_rules.csv} 中选取支持度、置信度和提升度较高的规则，解释“描述层级、干扰条件、临界性”如何共同诱发某类判断或错误模式。该部分是体现数据挖掘方法价值的重点。
